

% \paragraph{Section Overview}
% We summarize what the prototype Neural Computers (NCs) in this paper can reliably do today and identify the bottlenecks that prevent general-purpose use, highlighting where current world model~\citep{ha2018world,openai2024sora,polyak2024movie,deepmind_veo_2025} approaches fall short. We then contrast 
% % \wenyi{compare with TC, TC+Agent = GCC and CAU, and world models.} 
% NCs with conventional computers clarify what is uniquely enabled by learned neural systems. Finally, we define Completely Neural Computers (CNCs) and outline a roadmap toward them.

% \subsection{From Neural Computers to Completely Neural Computers}

% \paragraph{Current Status of NCs}
% Our CLI and GUI-based neural computers demonstrate that a learned latent runtime can implement components of the classical stack with measurable interface fidelity. In terminal environments, character-level OCR accuracy reaches $0.54$ (\Cref{tab:cligen-ocr}), while desktop cursor control achieves $98.7\%$ accuracy with explicit visual supervision (\Cref{tab:cursor-loss}). In GUIWorld, $110$ hours of goal-directed trajectories outperforms approximately $1{,}400$ hours of random exploration (\Cref{tab:gui-data-quality}), underscoring the role of aligned supervision. These results suggest that existing technology supports reliable I/O and short-horizon control: current NCs can manage low-level interface primitives with moderate training costs given high-quality datasets.

% However, to qualify as general-purpose systems, NCs must go beyond basic I/O and short-term execution.
% At a minimum, NCs must be Turing complete~\citep{turing1936computable,siegelmann1992computational}, support universal programmability~\citep{von1993first,wilkes1981best}, and exhibit consistent runtime behavior~\citep{queloz2025explainability}. 
% While certain sequential neural architectures are Turing complete~\citep{siegelmann1992computational,perez2021attention} in principle, turning a trained instance into a reliably programmable model remains challenging.
% Ensuring stable behavior over long temporal horizons remains an open problem in neural systems~\citep{kirkpatrick2017overcoming,calanzonelogically}. Section~\ref{sec:roadmap} provides a more detailed discussion on these requirements.
% To the best of our knowledge, existing works in world models lack a study of the computability class of the resulting learned models.


% \paragraph{Fundamental differences between NCs and conventional computers}
% We compare NCs and conventional computers. 
% More precisely, we refer to conventional computers as random-access machines~\citep{hartmanis1974power} 
% with instruction set architecture~\citep{anagnostopoulos1973computer} 
% that are programmable in human-designed high-level languages~\citep{backus1957fortran,backus1963revised}.
% % Their CUA and GCC extensions introduce agents as an intermediate layer between users and the I/O interfaces of conventional computers, so that users do not directly operate or program the computer but instead interact with it through these agents.
% NCs differ fundamentally from conventional computers in their architectural and programming-language semantics.

% \wenyi{
% Differences to conventional computers:
% 1. architectural semantics:
% conventional computers have local and symbolic semantics, making them precise but brittle. NCs on the other hand have holistic and numerical semantics, making them coarse but robust. Although NCs and conventional computers can possibly emulate each other, such emulation adds another layer of unnecessary difficulty when the nature of applied tasks are known.
% 2. learned v.s. human-designed programming-language semantics 
% Since NCs are learned from data, their programming-language semantics are not limited by human design. As an example, LLMs can be viewed as programmable computers with prompts being programs. In this view, the the programming language that is used to program an LLM can be a natural language, which no none-neural system has been successful developed to interpret it. Such learned programming-language semantics ...
% }

% \wenyi{
% An extension to conventional computers is  CUA. 
% }


% NCs and GCC diverge in their computational architecture. In GCC, a foundation model operates a classical computer through low-bandwidth I/O (e.g., mouse actions and screen frames), forcing both sides to communicate through compressed commands and observations. Each component only has direct access to its own internal state. In contrast, NCs unify the execution pipeline within a neural representation, which can reduce duplicated computation across components and lower synchronization overhead.

% \paragraph{The unique benefits of NCs and definition of CNC}
% A unified neural architecture can, in principle, expose interactions that are difficult in GCC. We group these potential benefits into three categories:

% \begin{enumerate}
% % \item \textbf{Learned computational semantics:}
%     \item \textbf{Shared computation:} Unlike GCC, which runs the foundation model and the classical computer separately, NCs unify computation, making resources available throughout the process. Typically, GCC operates where the classical computer has far less compute than the foundation model. An NC instance (fixed weights) could potentially be programmed to implement and execute new AI functions that require massive parallel computation.
%     \item \textbf{Shared information:} GCC relies on a low-bandwidth interface between the model and the computer. This bottleneck prevents the model from directly accessing large on-machine data (e.g., files), and prevents the classical runtime from accessing dense internal model states. NCs avoid this hard separation, allowing in-context-developed functions to communicate with the NC via learned semantics and exchange information at higher bandwidth.
%     \item \textbf{Application development via numerical optimization:} Because the memory state of an NC is numerical, one can frame some forms of ``program synthesis'' as optimizing that state under an objective. When the objective is differentiable through the model, first-order methods such as Adam \citep{adam2014method} apply; when it is not, gradient-free methods such as Natural Evolution Strategies \citep{wierstra2014natural} provide an alternative.
% \end{enumerate}

% Finally, we define a \textbf{Completely Neural Computer (CNC)} as a Neural Computer instance that is (1) Turing complete, (2) universally programmable, (3) consistent in its output unless explicitly reprogrammed, and (4) behaviorally exhibits the unique architectural benefits of shared computation, shared information, and a numerically optimizable memory state.

% \subsection{A Roadmap Towards CNC}
% \wenyi{add some recent techniques}
% \label{sec:roadmap}
% We frame the CNC transition as answering six scientific questions corresponding to its defining properties.

% \paragraph{Turing completeness}
% A Neural Computer (NC) instance (a specific architecture with fixed weights) defines a class of computational models in which each model corresponds to at least one memory state instance.
% To demonstrate that an NC instance is Turing complete, one must show it implements a Universal Turing Machine (UTM). In practice, this means that for any given Turing machine, there exists an initial memory state that allows the NC to emulate that machine exactly: it must produce identical outputs for halting inputs and execute indefinitely otherwise.

% For an NC instance to achieve universality, two elements are essential:

% \begin{enumerate}
%     \item  Iteration and Recursion: This is naturally facilitated by the recurrent architecture inherent to NCs.
% \item Unbounded Effective Memory: An NC instance has an unbounded effective memory if there are infinitely many possible memory state instances, each corresponding to a unique Turing machine.
% \end{enumerate}

% \paragraph{Universal programmability}
% An NC is universally programmable if, for each Turing machine, there exists an input sequence such that the NC realizes a new memory state representing that machine. This is likely to be achieved through compositional neural routines and partially persistent memory to store them. Progress can be tracked by routine reuse and compositional generalization.

% \paragraph{Runtime consistency}
% A CNC must preserve its function unless explicitly reprogrammed. For each memory state, there must be a non-empty set of inputs that executes the CNC without changing its pure function.
% Notice that since CNCs are Turing complete, for each possible input, there must be at least one memory state instance where this invariance does not hold. Thus, the set of inputs that preserve the CNC's function must be memory state dependent.
% We hypothesize that gating mechanisms, such as those in LSTM \citep{hochreiter1997long}, are effective in achieving such memory-state-conditioned invariance.

% \paragraph{Shared computation}
% CNCs leverage computational benefits to develop cognitive functions in context, closing the loop from neural execution to neural creation. Furthermore, similar to current AI's primitive self-improvement loop through recursive code generation, CNCs generalize this into a self-improvement cascade. CNCs emit reusable neural artifacts that can be used to generate future neural artifacts. Progress can be tracked by the reusability of cognitive functions emitted by the NC that were not present in its initial state and by whether self-improvement occurs.

% \paragraph{Shared information}
% In a GCC architecture, information transfer is restricted to high-level, human-interpretable interfaces: a literal ``keyhole'' of pixels and keystrokes. In contrast, the CNC architecture facilitates shared information by allowing neural processes to access the internal state of the environment and vice versa without serialization. A CNC achieves this when it can process and manipulate large-scale data structures (e.g., dense vector databases or raw telemetry at scale) without the latency or resolution loss associated with traditional I/O.

% Validation of shared information requires demonstrating that an in-context developed function can query the NC's latent memory space at a bandwidth exceeding that of a standard visual/textual interface. Progress can be measured by the latent-throughput ratio: the amount of relevant information exchanged between a task-specific routine and the system state per computational step, compared to the same task performed over a simulated CLI or GUI on classical computers.

% \paragraph{Application development via numerical optimization}
% The final milestone toward CNC realization is the transition from manual ``programming'' to gradient-based state synthesis. Because the memory of a CNC is a continuous numerical manifold, finding a memory state that performs a specific task can be framed as an optimization problem. A system exhibits this property when a user can define a loss function (e.g., ``minimize the error in this mathematical proof'') and use numerical solvers to directly update the NC's memory state to satisfy that objective.

% Unlike GCC, where optimization requires discrete code generation and a ``compile-and-test'' loop, the CNC allows for differentiable programming of the computer itself~\citep{innes2019differentiable}. Success is measured by the convergence efficiency of numerical solvers in finding valid program states compared to the success rate of combinatorial search (i.e., LLM-based code generation~\citep{hong2023metagpt}).
\paragraph{Section Overview}
We summarize what the prototype Neural Computers (NCs) in this paper can reliably do today and identify the bottlenecks that prevent general-purpose use, highlighting where current world model~\citep{ha2018world,openai2024sora,polyak2024movie,deepmind_veo_2025} approaches fall short.
% \mc{oh I see, I think this sentence is correct. My mistake in Wechat.} 
We then contrast NCs with conventional computers and define a class of general-purpose NCs, namely Completely Neural Computers (CNCs).
Finally, we outline a roadmap toward CNCs, discuss their relationship to other foundation models, and share additional thoughts on NCs.

\subsection{From Neural Computers to Completely Neural Computers}

\paragraph{Current Status of NCs}
Our CLI and GUI-based neural computers demonstrate that a learned latent runtime can implement components of the classical stack with measurable interface fidelity. In terminal environments, character-level OCR accuracy reaches $0.54$ (\Cref{tab:cligen-ocr}), while desktop cursor control achieves $98.7\%$ accuracy with explicit visual supervision (\Cref{tab:cursor-loss}). In GUIWorld, $110$ hours of goal-directed trajectories outperform approximately $1{,}400$ hours of random exploration (\Cref{tab:gui-data-quality}), underscoring the role of aligned supervision. These results suggest that existing technology supports reliable I/O and short-horizon control: current NCs can manage low-level interface primitives with moderate training costs given high-quality datasets.

However, to qualify as general-purpose systems, NCs must go beyond basic I/O and short-term execution. 
At a minimum, NCs must be Turing complete~\citep{turing1936computable,siegelmann1992computational}, support universal programmability~\citep{von1993first,wilkes1981best}, and exhibit consistent runtime behavior~\citep{queloz2025explainability}. 
% \mc{Here, could we illustrate more, like make `Turing complete`, `support universal programmability`, `exhibit consistent runtime behavior` in to items, and each follows by a short but correct sentence. Otherwise, these are just 3 fancy terms, and readers could not really get at once.} \mc{Okay, I found you have illustrate in 4.2}
While certain sequential neural architectures are Turing complete~\citep{siegelmann1992computational,perez2021attention} in principle, turning a trained instance into a reliably programmable model remains challenging.
Preliminary attempts, including Neural Virtual Machine~\citep{katz2019programmable} and NeuroLISP~\citep{davis2022neurolisp}, have been explored.
% For example, Neural Turing Machines (NTMs)~\citep{graves2014neural} and Differentiable Neural Computers (DNCs)~\citep{graves2016hybrid} are Turing complete in the asymptotic sense and have been shown to learn to follow specific algorithmic behavior. However, no ex
Furthermore, ensuring stable behavior over long temporal horizons remains an open problem in neural systems~\citep{kirkpatrick2017overcoming,calanzonelogically}. 
% \mc{Here is is better to mention Differential NC and Neural Turning Machine's advantages  and drawback.}
Section~\ref{sec:roadmap} provides a more detailed discussion of these requirements. 
% \mc{The following sentence give me a feeling, that we don't know: 1. what's the ideal world models, 2. what's the ideal NC? Actually, we utilize current "world model"-alike pipeline, but this is just the technique and it is imcomplete for CNC. So shoudld NC includes part of WM, or shoud WM includs part of NC, here need to be more clearly illsutrated...}
To the best of our knowledge, existing work on world models lacks an analysis of the computability class of the learned models.
See Section~\ref{sec:relations} for more discussion between NCs and other foundation models, including world models.

\paragraph{Fundamental differences between NCs and conventional computers}
We compare NCs and conventional computers.
More precisely, we refer to conventional computers as random-access machines~\citep{hartmanis1974power}
with instruction set architecture~\citep{anagnostopoulos1973computer}
that are programmable in human-designed high-level languages~\citep{backus1957fortran,backus1963revised}.
% Their CUA and CUA extensions introduce agents as an intermediate layer between users and the I/O interfaces of conventional computers, so that users do not directly operate or program the computer but instead interact with it through these agents.
NCs differ fundamentally from conventional computers in their architectural and programming-language semantics.

% \mc{I agree with you that miss citations...}
At the architectural level, random-access machines instantiate local, compositional symbolic semantics~\citep{Newell1976}, yielding exactness and interpretability, but brittleness under noise and model mismatch~\citep{BROOKS1991139}. 
Neural computers, by contrast, realize holistic, distributed numerical semantics, trading precise local semantics for robustness and generalization~\citep{Hinton1986,Bishop2006}. 
Empirical evidence indicates that such holistic numerical representations are particularly well suited to domains characterized by high-dimensional representations~\citep{bengio2013representation}, soft or statistical constraints~\citep{Smolensky1988-SMOOTP-2}, and globally coupled structures~\citep{silver2017mastering,vaswani2017attention}, including perception, natural language, motor control, planning under uncertainty, and approximate reasoning~\citep{lecun2015deep}.
Although conventional computers can, in principle, emulate NCs, doing so often introduces unnecessary conceptual and engineering complexity when the target tasks are already well matched to neural architectures.

At the programming-language level, NCs differ from conventional computers because their ``language semantics'' are the meanings of user input sequences learned from data rather than explicitly designed by humans. For example, LLMs can be viewed as programmable computers in which prompts act as programs~\citep{Reynolds2021}. In this case, the programming language is a natural language, which no non-neural system has historically been able to interpret robustly at scale~\citep{jm3}. 
More broadly, learned programming-language semantics are not constrained by a human-specified syntax/semantics boundary and can, therefore, encode task-relevant conventions implicitly~\citep{wei2023chainofthoughtpromptingelicitsreasoning}.


% \textbf{Application development via numerical optimization:} Because the memory state of an NC is numerical, one can frame some forms of ``program synthesis'' as optimizing that state under an objective. When the objective is differentiable through the model, first-order methods such as Adam \citep{adam2014method} apply; when it is not, gradient-free methods such as Natural Evolution Strategies \citep{wierstra2014natural} provide an alternative.

\paragraph{Definition of Completely Neural Computers} Thus, we define a Neural Computer instance to be complete if it is (1) Turing complete, (2) universally programmable, (3) consistent in its output unless explicitly reprogrammed, and (4) behaviorally realizes the advantages of the architectural and programming-language semantics of NCs relative to conventional computers.
We illustrate the aforementioned requirements in detail in the following section.

\subsection{A Roadmap Towards CNC}
% \textcolor{red}{add some recent techniques}
\label{sec:roadmap}
We frame the CNC transition as answering five scientific questions corresponding to its defining properties.

\paragraph{Turing completeness}
A Neural Computer (NC) instance (a specific architecture with fixed weights) defines a class of computational models in which each model corresponds to at least one memory state instance.
To demonstrate that an NC instance is Turing complete, one must show it implements a Universal Turing Machine (UTM). In practice, this means that for any given Turing machine, there exists an initial memory state that allows the NC to emulate that machine exactly: it must produce identical outputs for halting inputs and execute indefinitely otherwise. 

Notice that although Neural Turing Machines (NTM)~\citep{graves2014neural} and Differentiable Neural Computers (DNC)~\citep{graves2016hybrid} are Turing complete in the asymptotic sense, a particular NTM or DNC instance with finite precision cannot be Turing complete due to their fixed finite memory size.
For an NC instance to achieve universality, an unbounded effective memory is necessary.
% \begin{enumerate}
%   \item  Iteration and Recursion: This is naturally facilitated by the recurrent architecture inherent to NCs.
%   \item Unbounded Effective Memory: 
An NC instance has an unbounded effective memory if there are infinitely many possible memory state instances, each corresponding to a unique Turing machine. Existing works achieve such unboundedness by progressively growing model parameters~\citep{rusu2016progressive} or context~\citep{vaswani2017attention}.
% \end{enumerate}

\paragraph{Universal programmability}
An NC is universally programmable if, for each Turing machine, there exists an input sequence such that the NC realizes a new memory state representing that machine.
Most existing computational universality results for neural networks are established by constructing computational primitives and proving that their composition can simulate a universal computational model.
Likewise, we believe universal programmability in NCs can be achieved through compositional neural routines~\citep{reed2015neural}.


\paragraph{Runtime consistency}
A CNC must preserve its function unless explicitly reprogrammed. For each memory state, there must be a non-empty set of inputs that executes the CNC without changing its pure function.
% Notice that since CNCs are Turing complete, for each possible input, there must be at least one memory state instance where this invariance does not hold. Thus, the set of inputs that preserve the CNC's function must be memory state dependent.
Operationally, this likely benefits from a reliable programming interface that induces intentional functional updates and an execution interface in which ordinary inputs predominantly use the programmed state rather than modify it. 
This motivates training and architectural mechanisms that disentangle function updates from function use, so that users can install and compose routines without accidental drift.
We hypothesize that gating mechanisms, such as those in LSTM \citep{hochreiter1997long}, are effective in achieving this conditional invariance.

\paragraph{Architectural semantics}
Since NC behavior is governed by real-valued parameters, learning can produce novel input–output mappings beyond those seen in training~\citep{poggio2019theoretical}, potentially introducing new primitives~\citep{ha2016hypernetworks}.
The combination of such new primitives could enable qualitatively new functions, yielding out-of-distribution functional generalization~\citep{lake2018generalization}.
% Progress should therefore be evaluated by examining NCs on held-out function classes absent from training, using metrics such as exact output matching and stability under small input perturbations.
Beyond emulating conventional computers, NCs can natively support AI-native functions whose semantics are ill-suited to symbolic APIs~\citep{marcus2018deep}, including probabilistic inference over high-dimensional latent states~\citep{kingma2013auto}, representation learning~\citep{bengio2013representation}, retrieval over dense memories~\citep{graves2016hybrid}, and end-to-end differentiable pipelines that couple perception and control~\citep{silver2017mastering}. These functions are first-class at the architectural level and operate directly on distributed states~\citep{lecun2015deep}. This enables capabilities such as learned heuristics~\citep{silver2017mastering}, uncertainty-aware decision-making~\citep{clements2019estimating}, and continual adaptation~\citep{parisi2019continual}.

Because the memory state of an NC/CNC is numerical, computer configuration and design emerge as alternatives to application-level programming: the computer itself is configured by optimizing its internal state to achieve desired computational behaviors under task-defined objectives. 
Depending on the differentiability of the loss, methods such as Adam \citep{adam2014method} and \citep{wierstra2014natural} apply. In a CNC, the memory constitutes a continuous manifold, so realizing a target capability amounts to synthesizing a machine configuration (a memory state) that minimizes a user-specified loss (e.g., “minimize proof error”) via direct numerical updates to the computer’s state. This reframes system construction from discrete code authoring to differentiable configuration of the computer itself \citep{innes2019differentiable}, with progress evaluated by solver convergence, stability, and reliability relative to combinatorial program search (e.g., LLM-based code generation \citep{hong2023metagpt}).

\paragraph{Programming-language semantics}
The learned programming-language semantics of NCs enable a shift from rigid coding to learned specifications, in which user inputs themselves function as programs~\citep{brown2020language,wei2022chain}. 
Rather than centering development on explicitly authored code, NCs expose a learned language whose syntax and semantics are acquired from data~\citep{radford2019language,bommasani2021opportunities}, so natural-language instructions, examples, and constraints serve as executable specifications~\citep{ouyang2022training}. 
Scrupulous computation is internalized within the model’s learned substrate and steered by contextual conditioning.
Consequently, brief user inputs can replace long sequences of low-level actions. 
Development, therefore, moves from code authoring to curating, specifying, and verifying inputs under a learned programming-language semantics, aligning system behavior with human intent via in-context specification rather than forcing users to conform to rigid, brittle interfaces~\citep{austin2021program}.

Since NCs are programmed via users' input sequences under learned programming-language semantics, the training data for programming NCs, i.e., paired user I/O traces~\citep{flener2008introduction}, is far more abundant and continuously generated than high-quality, human-written code.
Every interaction with digital systems produces structured streams of inputs, interface states, and effects that can be logged at scale (e.g., keystrokes, cursor trajectories, screen transitions), yielding orders-of-magnitude more supervision than curated program corpora~\citep{yao2022webshop}. 
These I/O traces constitute executable specifications, revealing user intentions and computer behavior~\citep{cypher1993watch}.
This enables end-to-end learning of interface conventions, control policies, and task semantics without requiring explicit program text~\citep{yao2022react}. This asymmetry in data availability favors NC training regimes that leverage ubiquitous interaction logs and, by supporting broader task coverage, reduces dependence on brittle, sparsely available code datasets.

\subsection{Relations to Other Foundation Models}
\label{sec:relations}
\paragraph{World Models}
World models learn environment dynamics by predicting action-conditioned transitions~\citep{ha2018world}.
Such target environments range from the most ambitious—the whole world~\citep{schmidhuber1990making}—to much narrower scopes, such as a robot arm~\citep{feng2023finetuning}.
From a world model perspective, neural computers are world models that focus on a specific type of environment, i.e., interactive computers.
Interestingly, many current approaches to modeling computational environments, such as the physical world, heavily rely on computer-generated data~\citep{richter2016playing,tobin2017domain,dosovitskiy2017carla,garcia2023robust}, potentially leading to models that share characteristics with neural computers. 
\mc{Better to add an image of how NC distinguished or related to WM in different situations}

\paragraph{Computer Use Agents}
Another interesting approach that integrates foundation models (typically LLMs) and conventional computers is the computer use agent (CUA)~\citep{openaiComputerUsingAgent,claudeComputerTool,sager2025comprehensivesurveyagentscomputer}. 
CUA introduces LLM-powered agentic systems as an intermediate layer between users and the I/O interfaces of conventional computers, so that users do not directly operate or program the computer but instead interact with it through these agents.
This provides a computational paradigm that leverages both the cognitive ability of LLMs and the preciseness and programmability of conventional computers~\citep{song2025coact1computerusingagentscoding}.
However, this paradigm is constrained by the low-bandwidth I/O interface and the strict separation between the LLM and the conventional computer, which can introduce latency and restrict access to internal states.
Moreover, application development with CUAs still typically occurs on the conventional computer under fixed, human-designed programming-language semantics, which encourages brittle symbolic programs~\citep{liang-etal-2017-neural} and limits end-to-end optimization over richer internal representations~\citep{hu2017learningreasonendtoendmodule,gaunt2016terpretprobabilisticprogramminglanguage}.
We hypothesize that a sufficiently capable NC can internalize these ``agentic functions'' within its internal neural system.

\subsection{Additional Thoughts} %Inspired by This Research}

\mc{technological equality}

\mc{out of hardware limitation}

\mc{NC generate NN}

\mc{universal tensor file representation}

\mc{extremely sparse models}

\paragraph{Video models as a pragmatic prototype substrate}
We build our prototypes on state-of-the-art video models because they currently provide the simplest path to an end-to-end learned latent runtime that jointly models pixels, dynamics, and action-conditioned control. This choice is pragmatic rather than fundamental. In our experiments, symbolic and algorithmic reasoning in terminal settings remains inconsistent for most strong video models, and even simple arithmetic can fail (\Cref{tab:cligen-arith}). Sora2 is a notable exception in our probe, achieving $71\%$ arithmetic accuracy, suggesting that some terminal symbolic reasoning is already possible in modern video generators. At the same time, we do not claim that video models cannot reason more broadly: recent work reports that video models can act as zero-shot learners and reasoners in naturalistic settings \citep{wiedemer2025video}. We expect reasoning capabilities to improve quickly with continued progress in video modeling, but our results suggest that CNC-level reliability will likely require additional architectural and training ingredients beyond scaling today's video generators.

\paragraph{A hypothesis: machine-native neural architectures}
We emphasize that the following is a conjecture rather than a conclusion drawn from our experiments. Closing the reasoning gap may not require designing neural networks that more closely mimic animal cognition or the human brain.
Many influential architectures, including convolutional networks \citep{fukushima1980neocognitron} and Transformers \citep{vaswani2017attention}, are highly engineered systems, but their core inductive biases remain strongly influenced by biological perception and attention. These models primarily rely on continuous, distributed representations, in which reasoning behavior emerges implicitly from large-scale training.
We hypothesize that CNCs may instead benefit from designs that are explicitly machine-native. Developing discrete operations, compositional structures, and verifiable computation that are harmonious in neural systems may play an essential role in designing such systems. This approach follows more closely the construction of classical computers from well-defined computational primitives and stands in contrast to relying on emergent reasoning in generic video generation models.
